% Independent mathematical derivation, 2026-09-23.
% DOR1--22. DOR20--22 are the root's exact CCM receiving application.
% All tensor products and coinvariants are algebraic.
\section{The odd source sector and the full range of the diagonal receiver}
\label{sec:diagonal-odd-sector-range}

\paragraph{Sources and exact scope.}
This calculation continues DMR1--29 in
\nolinkurl{DIAGONAL_MOMENT_REFINED_RECEIVER.tex}, read in full for this
derivation. The original Schwartz source and rational periodization are
those of Alain Connes,
\href{https://arxiv.org/abs/math/9811068v1}{\emph{Trace formula in
noncommutative geometry and the zeros of the Riemann zeta function}},
Section III, and Alain Connes, Caterina Consani and Matilde Marcolli,
\href{https://arxiv.org/abs/math/0703392v1}{\emph{The Weil proof and the
geometry of the adeles class space}}, source labels
\texttt{rangecVdef}, \texttt{varthetaH1V}, \texttt{vanishV},
\texttt{tracepairing}. Their use here is through the precise original
definitions retained in DMR1 and the ABR source-coordinate comparison;
no new full-paper reading is claimed. The full support carrier has the
originating byline The Clankers,
\href{https://github.com/KokunoYumeto/zeta-function-research-reader/blob/a2851f9eb352e7e4376bc629de86fa4ed9c1c434/workbenches/splitzero-tandem/foundations/split-support-geometry-v11/split_support_geometry_arithmetic_curve_v11.tex}{\emph{Split Support Geometry}, v11,
Definition \texttt{def:lattice-split}}.

We compute the previously undetermined $H_{00}$ range of DMR18,
prove the full receiver onto, and give explicit coordinates for an
entire source summand and its kernel. The real odd source is retained.
It is killed by independent spherical averaging but survives the
original diagonal multiplication. No assertion of Weil positivity,
purity or an RH conclusion is made.

\subsection{An exact odd summand of the original coefficient kernel}

Use the original spaces, measures and actions
\begin{equation}\tag{DOR1}
\begin{gathered}
T=\mathcal S(\mathbb R)\otimes\mathcal S(\mathbb A_f),\qquad
G=\mathbb Q^\times,\qquad (U_a\xi)(x)=\xi(a^{-1}x),\\
E\xi(x)=\sum_{q\in\mathbb Q^\times}\xi(qx),\qquad
K=\ker(b\otimes b)\cap\ker(E\otimes E),\\
(\mu_yt)(z)=t((y_1,1_f)z,(y_2,1_f)z),\qquad
m_y(k)=\int_{\mathbb A}\mu_yk,\qquad y_1,y_2>0.
\end{gathered}
\end{equation}
Here $b$ is the complete tensor of local evaluation and local integral
maps, not just its two constant endpoints. Lebesgue measure is used
at infinity and $\operatorname{vol}(\mathbb Z_p)=1$ at every prime.
The action on $T\otimes T$ is diagonal. Every tensor product, direct
sum and coinvariant below is algebraic; every vector has finite tensor
length. No topological closure of a relation space is introduced.

Let $P_f$ average finite units and let $J$ change only the real sign:
\begin{equation}\tag{DOR2}
\begin{gathered}
P_f\xi=\int_{\widehat{\mathbb Z}^{\times}}U_{(1,u)}\xi\,du,
\quad \int du=1,\qquad J=U_{(-1,1_f)},\\
P^-=\tfrac12(I-J)P_f,\quad
T^-=P^-T,\quad H^-=\mathcal S_{\mathrm{odd}}(\mathbb R),
\quad D=T^-\otimes T^-.
\end{gathered}
\end{equation}
The finite average is finite on each original finite Bruhat test:
its compact unit orbit factors through a finite quotient of
$\widehat{\mathbb Z}^{\times}$. Thus DOR2 is an actual map on $T$.
It is idempotent and commutes with every idele dilation. An element
of $T^-$ is odd in the real variable and invariant under every finite
unit. Its local real evaluation and integral both vanish. Hence
$b(T^-)=0$, with every local support pattern retained.

Finite unit invariance also gives $\xi(-x)=-\xi(x)$ for $\xi\in T^-$,
where this minus sign changes every adelic coordinate. The original
periodization on an idelic $x$ is absolutely convergent: the finite
support forces its rational summation parameters into a fixed
fractional integer lattice, and real Schwartz decay gives a convergent
series there. Pairing each $q$ with $-q$ therefore proves
\begin{equation}\tag{DOR3}
E(T^-)=0,\qquad D\subset K,\qquad
P^-\otimes P^-:K\longrightarrow D\text{ is an equivariant retraction}.
\end{equation}
The retraction assertion does not presume that independent averaging
commutes with $\mu_y$. Its image is contained in $D\subset K$, and
its restriction to $D$ is the identity. Therefore the inclusion
$H_0(G,D)\to H_0(G,K)$ is injective, with the induced retraction as
its left inverse. This is an actual source summand.

Put $R=\mathbb C[\mathbb Q_{>0}^{\times}]$. The map
\begin{equation}\tag{DOR4}
\Theta_-(t_a\otimes F)(x)=F(x_\infty/a)
                              \mathbf1_{a\widehat{\mathbb Z}}(x_f)
\end{equation}
is an isomorphism $R\otimes H^-\to T^-$. To verify its finite part,
a compactly supported locally constant radial function on
$\mathbb Q_p$ is a finite linear combination of the nested ball
indicators $\mathbf1_{p^n\mathbb Z_p}$. Such finite combinations
are independent by evaluating on the successive nonempty shells.
At finitely many primes take their tensor products; at all remaining
primes keep $\mathbf1_{\mathbb Z_p}$. The resulting finite ball
products are indexed by the positive rational number $a$ with the
specified valuations. These products form a basis of the finite
unit invariant Bruhat space. Rescaling each real coefficient by $a$
then gives DOR4 and its unique inverse. Positive rational $q$ acts
by $t_a\mapsto t_{qa}$, and rational $-1$ acts as $-I$.

On $D$ the rational sign is therefore $+I$. Dividing the pair of
positive rational indices $(a,b)$ by their common second index
identifies its diagonal $R$-module as free. Consequently
\begin{equation}\tag{DOR5}
\begin{gathered}
H_0(G,D)\cong
\mathcal A^-:=\bigoplus_{c\in\mathbb Q_{>0}^{\times}}
                         (H^-\otimes H^-)_c,\\
[\Theta_-(t_a\otimes F)\otimes\Theta_-(t_b\otimes G)]
\longmapsto [a/b]\otimes F\otimes G.
\end{gathered}
\end{equation}
An inverse sends $[c]\otimes F\otimes G$ to the class with
indices $(c,1)$. Indeed applying the common rational $b^{-1}$
changes $(a,b)$ to $(a/b,1)$ without changing $F$ or $G$.
There are no further relations: the common index is precisely a
free $R$ factor, and its coinvariant is the augmentation of that
factor. This proves both inverse identities, including for sums.

\subsection{The complete multiplication receiver on this summand}

Write $c=m/n$ in its unique relatively prime positive integer form.
This is only the original prime-valuation description of the rational
index: no source coefficient is divided out. For the representative
$(c,1)$, intersection of the two finite balls is
$c\widehat{\mathbb Z}\cap\widehat{\mathbb Z}
=m\widehat{\mathbb Z}$. The pullback and its actual ABR coefficient
are therefore
\begin{equation}\tag{DOR6}
\begin{split}
\mu_y\bigl(\Theta_-(t_{m/n}\otimes F)
                      \otimes\Theta_-(t_1\otimes G)\bigr)(z)
 &=F(y_1nz_\infty/m)G(y_2z_\infty)
                                  \mathbf1_{m\widehat{\mathbb Z}}(z_f)\\
 &=\Theta\bigl(t_m\otimes H_{m,n}\bigr)(z),\\
H_{m,n}(x)&=F(y_1nx)G(y_2mx).
\end{split}
\end{equation}
The finite volume is $1/m$, and the real substitution $z_\infty=mx$
gives the exact moment $\int_{\mathbb R}H_{m,n}(x)\,dx$.
The result is already invariant under finite units and the real
sign, so $Q\mu_y=\mu_y$ on $D$.

Define the linear map on the explicit direct sum
\begin{equation}\tag{DOR7}
\begin{gathered}
\mathcal T_y:\mathcal A^-\longrightarrow
 H_{[0]}:=\{h\in\mathcal S_{\mathrm{even}}(\mathbb R):h(0)=0\},\\
\mathcal T_y\left(\sum_{m,n,j}[m/n]\otimes F_{m,n,j}\otimes G_{m,n,j}\right)(x)
 =\sum_{m,n,j}F_{m,n,j}(y_1nx)G_{m,n,j}(y_2mx),\\
\nu_y=\int_{\mathbb R}\mathcal T_y,\qquad
\mathcal A_y^{-,0}=\ker\nu_y.
\end{gathered}
\end{equation}
Every displayed sum is finite; each term retains both original real
scales and its reduced positive integer pair. Equations DOR5--6
prove that $\nu_y$ is the restriction of the original descended
moment $\overline m_y$, and that $\mathcal A_y^{-,0}$ embeds in
$S_y=\ker\overline m_y\subset H_0(G,K)$.

For a moment-zero class the ABR formulas of DMR10 on the chosen
representative give all three coordinates, not only the augmented
function:
\begin{equation}\tag{DOR8}
\begin{gathered}
\Phi=\sum_{m,n,j}H_{m,n,j},\qquad \beta_0=0,\\
\beta_1=\sum_p\ell_p\sum_{m,n,j}v_p(m)
                              \int_{\mathbb R}H_{m,n,j}(x)\,dx.
\end{gathered}
\end{equation}
The second identity follows because every real coefficient is the
product of two odd functions and has evaluation zero. The third
retains every prime valuation and integral. Changing the coefficient
lift may alter this third coordinate by the complete connecting
image DMR15; it does not alter either of the first two. The DMR
decomposition retains that integral-prime coordinate separately.
Accordingly its canonical receiver on this source summand is
\begin{equation}\tag{DOR9}
\widehat F_y|_{\mathcal A_y^{-,0}}=(\mathcal T_y,0)
  :\mathcal A_y^{-,0}\longrightarrow H_{00}\oplus j_0(W_{\rm pr}),
\qquad
H_{00}=\{h\in H_{[0]}:\int h=0\}.
\end{equation}
Here $\widehat F_y$ is exactly DMR18, with its separately retained
$j_1(W_{\rm pr})$ contribution. No replacement receiver has been
introduced.

\subsection{A proved Schwartz factorization with its construction}

We now prove the factorization needed to compute the range, rather
than assuming that an arbitrary Schwartz function has appropriate
Schwartz factors. For every even $h\in\mathcal S(\mathbb R)$ there
exists a positive even $A\in\mathcal S(\mathbb R)$ such that
\begin{equation}\tag{DOR10}
h/A^r\in\mathcal S(\mathbb R)\quad\hbox{for every integer }r\ge1.
\end{equation}
Choose increasing real numbers $T_N$, $N\ge1$, satisfying
$T_1\ge2$, $T_{N+1}\ge T_N+2$, and
\[
|h^{(j)}(x)|\le(1+x^2)^{-N^2}
\quad(0\le j\le N,\ \log(1+x^2)\ge T_N).
\]
Such numbers exist from the defining Schwartz bounds, with finitely
many derivatives specified at each step. Choose a smooth increasing
function $\chi$ equal to zero on $(-\infty,0]$ and one on
$[1,\infty)$; for definiteness it may be obtained on $(0,1)$ by
integrating $e^{-1/(u(1-u))}$ and dividing by its positive integral
over $(0,1)$, retaining that integral as its defining constant.
Define, for $t\ge0$,
\begin{equation}\tag{DOR11}
\phi(t)=\int_0^t\left(1+\sum_{N\ge1}\chi(u-T_N)\right)du,
\qquad A(x)=\exp\{-\phi(\log(1+x^2))\}.
\end{equation}
The sum is locally finite. For $T_N\le t<T_{N+1}$ one has
$0\le\phi(t)\le(N+1)t$. For every fixed integer $M$ its derivative
is at least $M$ for all sufficiently large $t$; integrating this
inequality gives $\phi(t)/t\to\infty$. Derivatives of $\phi$ of
order at least two are uniformly bounded, because their transition
intervals $[T_N,T_N+1]$ are disjoint. Its first derivative is at
most $t+2$, using $T_N\ge N+1$.

Repeated differentiation of $\phi(\log(1+x^2))$ is therefore
bounded by a polynomial in $1+|x|$ in every fixed derivative
order. Every derivative of $A$ is $A$ times a polynomial in those
derivatives. Since $\phi(t)/t\to\infty$, $A$ and all its
derivatives decay faster than every power of $1+|x|$. This proves
$A\in\mathcal S(\mathbb R)$, with strict positivity and evenness.
For $t=\log(1+x^2)\in[T_N,T_{N+1})$, $j\le N$, the bounds give
\[
|h^{(j)}(x)|/A(x)^r
\le\exp\{-[N^2-r(N+1)]t\}.
\]
For fixed $r,j$ the exponent tends to negative infinity faster
than any prescribed fixed negative multiple of $t$. Derivatives
of $A^{-r}$ are $A^{-r}$ times polynomially bounded factors by
the same chain-rule calculation. The product rule now proves every
Schwartz seminorm of $h/A^r$ finite. This proves DOR10 completely.
No infinite tensor or infinite sum of source vectors was used;
DOR11 constructs a single permissible real Schwartz function.

When $h(0)=0$ and $h$ is even, set
\begin{equation}\tag{DOR12}
G_h(x)=xA(x),\qquad F_h(x)=\frac{h(x)}{xA(x)},\qquad
F_h(0)=0,\qquad h=F_hG_h.
\end{equation}
Both functions are odd Schwartz functions. At zero this follows
from the exact Taylor formula
$h(x)=x^2\int_0^1(1-u)h''(ux)du$; its quotient by $xA(x)$
extends smoothly and is odd. Away from zero the division by $x$
and DOR10 imply all Schwartz bounds. This proves the needed
factorization with an explicit construction and handles the
original zero to its actual order. It does not assert a fixed
factor $A$ valid for every $h$, a linear choice of factors, or
an equivariant section of the factorization.

\subsection{The exact range and a concrete kernel}

For $h\in H_{[0]}$ put
\begin{equation}\tag{DOR13}
\begin{gathered}
f_{h,y}(x)=F_h(x_\infty/y_1)\mathbf1_{\widehat{\mathbb Z}}(x_f),
\qquad
g_{h,y}(x)=G_h(x_\infty/y_2)\mathbf1_{\widehat{\mathbb Z}}(x_f),\\
k_{h,y}=f_{h,y}\otimes g_{h,y}\in D\subset K,\qquad
\mu_y k_{h,y}=h(x_\infty)\mathbf1_{\widehat{\mathbb Z}}(x_f),\qquad
m_y(k_{h,y})=\int_{\mathbb R}h.
\end{gathered}
\end{equation}
Both original individual theta sums vanish by DOR3, not by an
assumed arithmetic relation. The displayed pullback proves
$\mathcal T_y$ onto $H_{[0]}$. For $h\in H_{00}$ the source lies
in $K_y^0$ and has exact ABR coordinates $(h,0,0)$. Thus
\begin{equation}\tag{DOR14}
\begin{gathered}
\widehat F_y([k_{h,y}]_K)=(h,0),\qquad
0\longrightarrow\mathcal N_y^-
 \longrightarrow\mathcal A_y^{-,0}
 \xrightarrow{\mathcal T_y}H_{00}\longrightarrow0,\\
\mathcal N_y^-=
\left\{\sum_{m,n,j}[m/n]\otimes F_{m,n,j}\otimes G_{m,n,j}:
\sum_{m,n,j}F_{m,n,j}(y_1nx)G_{m,n,j}(y_2mx)=0
\text{ for every }x\in\mathbb R\right\}.
\end{gathered}
\end{equation}
This is a kernel on the explicitly determined algebraic direct
sum DOR5, not an unspecified cohomological kernel. Its moment
condition is automatic from the displayed pointwise identity.
It is nonzero. Choose linearly independent odd Schwartz functions
$F,G$; the element in its $[1]$ summand
\[
F(\cdot/y_1)\otimes G(\cdot/y_2)
 -G(\cdot/y_1)\otimes F(\cdot/y_2)
\]
has zero image and is nonzero. Nonzero follows by applying an
algebraic linear functional on the span of the first two independent
factors that takes values $1,0$. Its image is then the nonzero
second factor $G(\cdot/y_2)$. The injection DOR3 means this remains
a nonzero class in $H_0(G,K)$; it is not an artifact of a chosen
presentation of the larger kernel.

The full refined receiver is onto:
\begin{equation}\tag{DOR15}
\widehat F_y:S_y\twoheadrightarrow H_{00}\oplus j_0(W_{\rm pr}),
\qquad
F_y:H_0(G,K_y^0)\twoheadrightarrow Q_0.
\end{equation}
Here is the exact surjectivity argument including the simultaneous
coordinates. DMR24--27 supplies for each prime a class $s_p\in S_y$
with $\beta_0\widehat F_y(s_p)=-\ell_p$. Its first component is an
actual $h_p\in H_{00}$. By DOR14 choose $s'_p$ with image $(h_p,0)$.
Then $-s_p+s'_p$ has image $(0,\ell_p)$. Finite combinations give
every evaluation-prime vector, and DOR14 gives every $H_{00}$
vector with zero evaluation-prime component. This proves the first
surjection. For the second use DMR15,
$F_y\delta_y=j_1$, together with DMR17--18; the missing integral-prime
component is reached with its original connecting map. This does
not identify the full kernel of $\widehat F_y$ with the smaller
explicit kernel $\mathcal N_y^-$. It computes the entire range and
computes the exact kernel on the displayed source summand.

For completeness, the whole integral-prime component can also be
reached from the odd summand itself. Put $F(x)=G(x)=xe^{-\pi x^2}$
and use the same real rescalings and finite references as DOR13.
Its tensor $v\in D$ has
\begin{equation}\tag{DOR16}
m_y(v)=\int_{\mathbb R}x^2e^{-2\pi x^2}dx
       =\frac1{2\pi\,2^{3/2}}>0.
\end{equation}
The Gaussian identity follows by differentiating
$\int e^{-\pi A x^2}dx=A^{-1/2}$ at $A=2$, retaining its full
factor. For each prime $p$, the vector
$(U_p\otimes U_p-I)v/m_y(v)$ belongs to $D\cap K_y^0$ and has
ABR coordinates $(0,0,\ell_p)$. All sources in $D\cap K_y^0$
have $\beta_0=0$, by expansion in the original finite-ball basis
and odd real coefficients. Consequently
\[
F_y\bigl(\operatorname{im}[H_0(G,D\cap K_y^0)
                  \to H_0(G,K_y^0)]\bigr)
=H_{00}\oplus j_1(W_{\rm pr}).
\]
The integral-prime classes here may vanish after passage to
$H_0(G,K)$, exactly as DMR16 requires. They are retained in the
refined source and its DMR17 decomposition.

\subsection{Idele actions, theta image and full support}

For an idele $a$ retain its finite valuation factor
$r=\prod_p p^{v_p(a_p)}>0$ and original modulus
$b=|a_\infty|_\infty/r=|a|$. On $T^-$, finite units act trivially
and real sign $\sigma=\operatorname{sgn}(a_\infty)$ acts by
$\sigma$. Thus the exact action on DOR4 is
\[
U_a\Theta_-(t_c\otimes F)
 =\sigma\,\Theta_-(t_{rc}\otimes R_bF),\qquad
(R_bF)(x)=F(x/b).
\]
On $D$ the two signs multiply to $1$, the common finite rational
factor disappears by the original diagonal rational relation, and
the relative index remains unchanged. Therefore
\begin{equation}\tag{DOR17}
\begin{gathered}
U_a([c]\otimes F\otimes G)
 =[c]\otimes R_{|a|}F\otimes R_{|a|}G,\\
\mathcal T_yU_a=R_{|a|}\mathcal T_y,\qquad
\nu_yU_a=|a|\nu_y.
\end{gathered}
\end{equation}
This proves equivariance of DOR14 for the full original idele
action and proves that its kernel is invariant. DOR15 is
equivariant by the already proved DMR action; no equivariant
section of either surjection is claimed.

For the original source $k_{h,y}$ the independent spherical
projection and the diagonal theta receiver give, simultaneously,
\begin{equation}\tag{DOR18}
\begin{gathered}
(Q\otimes Q)k_{h,y}=0,\qquad
Q\mu_yk_{h,y}=h\mathbf1_{\widehat{\mathbb Z}},\qquad
(E\otimes E)k_{h,y}=0,\\
E\mu_yk_{h,y}(u,1_f)=2\sum_{n\ge1}h(nu),\qquad u>0,\\
\int_0^\infty E\mu_yk_{h,y}(u,1_f)u^{s-1}du
 =2\zeta(s)\int_0^\infty h(x)x^{s-1}dx,
\qquad\operatorname{Re}s>1.
\end{gathered}
\end{equation}
The first identity follows because each factor is odd under the
real sign average. The second is the actual product of two odd
functions and hence is even; it is not obtained by asserting
commutation of the two averages. The theta series retains the
two rational signs, the original zeta function and the original
Mellin measure. Its Mellin identity follows by absolute Fubini:
the integral of the absolute majorant is
$2\sum_{n\ge1}n^{-\operatorname{Re}s}
\int_0^\infty|h(x)|x^{\operatorname{Re}s-1}dx<\infty$.
As $h$ runs over $H_{00}$, this is the entire source theta image
in DMR28, not an uncomputed subset of it. Its further image in
the arithmetic quotient by $\overline{E(T_0)}$ remains zero,
because these are actual elements of $E(T_0)$. This calculates
the surviving source range without converting it into an
unsupported spectral assertion.

Finally, for the original bounded distributive support lattice $L$
retain all labels, their original joins and meets, and
\begin{equation}\tag{DOR19}
\begin{gathered}
G_L(M)=\{(0,\lambda):\lambda\in L\}\cup(M\times\{1_L\}),
\quad e=(0,1_L),\quad\tau=(0,0_L),\\
A^\uparrow(v,\lambda)=(Av,\lambda)
\quad\text{for every displayed linear map }A.
\end{gathered}
\end{equation}
The quotient uses exactly the same-label amplitude congruence
of DMR29. DOR3, DOR5, DOR7, DOR14 and DOR15 therefore have
these common-label lifts. In particular the construction of
$k_{h,y}$ is a construction of an amplitude preimage, with its
one unchanged source label. It does not give two independently
assignable tensor labels. Odd cancellation in the original
individual theta factors sends top-supported amplitude to $e$,
not $\tau$, and the nonzero diagonal product retains that label.
The exact source distinction in DOR18 persists on the full
carrier. No ordinary additive homology on the label lattice
has been assumed.

\subsection{The original CCM degree adjustment inside this source summand}

We now apply the range construction to the actual function in
Connes--Consani--Marcolli, source lemma \texttt{degmodifyV}, with its
complete constants already reconstructed in ETC1--26 of the accompanying
\nolinkurl{CC_EXACT_TRIVIAL_CORRESPONDENCE.tex}. The endpoint values
and their two-dimensional inverse are existing results, not new claims
of this section. The new receiving calculation locates their complete
original family inside the odd tensor summand of the joint kernel.

For each original positive width $b$ define
\begin{equation}\tag{DOR20}
\begin{gathered}
\eta_b(x)=
\left(\pi^2(x/b)^4-\frac{3\pi}{2}(x/b)^2\right)e^{-\pi(x/b)^2},\\
F_b(x)=\pi(x/b)\left(\pi(x/b)^2-\frac32\right)
                  e^{-\pi(x/b)^2/2},\qquad
G_b(x)=(x/b)e^{-\pi(x/b)^2/2},\\
f_{b,y}(z)=F_b(z_\infty/y_1)\mathbf1_{\widehat{\mathbb Z}}(z_f),
\quad
g_{b,y}(z)=G_b(z_\infty/y_2)\mathbf1_{\widehat{\mathbb Z}}(z_f),
\quad t_{b,y}=f_{b,y}\otimes g_{b,y}.
\end{gathered}
\end{equation}
Multiplication gives $F_bG_b=\eta_b$ with every Gaussian exponent,
polynomial coefficient and real scale retained. Both factors are odd
Schwartz functions. Thus DOR3 and direct integration give
\[
t_{b,y}\in D\cap K_y^0,\quad (E\otimes E)t_{b,y}=0,\quad
(Q\otimes Q)t_{b,y}=0,\quad
\mu_y t_{b,y}=\eta_b\mathbf1_{\widehat{\mathbb Z}}.
\]
Indeed $\eta_b(0)=0$ and its additive integral is
$b(\pi^2\,3/(4\pi^2)-(3\pi/2)/(2\pi))=0$.
Its actual diagonal theta image is the original $f_b$ of ETC6:
$f_b(u)=2\sum_{n\ge1}\eta_b(nu)$.

The complete original finite data are an arbitrary finitely supported
integer vector $k=(k_p)_p$ and an independent real scale $a>0$.
Retain $N(k)=\prod_p p^{k_p}>0$, allowing negative $k_p$, and put
\begin{equation}\tag{DOR20a}
\begin{gathered}
f_{k,a,y}(z)=F_a(z_\infty/y_1)
                \prod_p\mathbf1_{p^{k_p}\mathbb Z_p}(z_p),\qquad
g_{k,a,y}(z)=G_a(z_\infty/y_2)
                \prod_p\mathbf1_{p^{k_p}\mathbb Z_p}(z_p),\\
t_{k,a,y}=f_{k,a,y}\otimes g_{k,a,y},\qquad b=a/N(k),\\
\mu_y t_{k,a,y}=\xi_{k,a}
 =\eta_a(z_\infty)\prod_p\mathbf1_{p^{k_p}\mathbb Z_p}(z_p),\\
t_{k,a,y}=(U_{N(k)}\otimes U_{N(k)})t_{b,y},\qquad
\int_{\mathbb A_\mathbb Q}\xi_{k,a}
=\frac{a}{N(k)}
 \left(\frac{\pi^2\,3}{4\pi^2}-\frac{(3\pi/2)}{2\pi}\right)=0.
\end{gathered}
\end{equation}
The finite product is exactly $\mathbf1_{N(k)\widehat{\mathbb Z}}$;
its square is itself and its original additive volume is
$\prod_p p^{-k_p}=N(k)^{-1}$. Both real factors remain odd,
so $t_{k,a,y}\in D\cap K_y^0$ and its independent theta and
independent $Q$ images vanish. The displayed rational comparison
has the explicit inverse $U_{1/N(k)}\otimes U_{1/N(k)}$.
At a finite-unit idele, a rational $q$ contributes exactly when
$v_p(q)\ge k_p$ for every prime, equivalently $q=N(k)n$ with
$n\in\mathbb Z\setminus\{0\}$. Consequently the diagonal theta
image is $f_{a/N(k)}$. The tuple $(k,a)$ remains in the antecedent;
the ratio $b$ does not recover it. In ABR coordinates
$\xi_{k,a}=\Theta(t_{N(k)}\otimes\eta_b)$, and its prime component
is exactly $\sum_p k_p\ell_p\otimes(\eta_b(0),\int\eta_b)
=\sum_p k_p\ell_p\otimes(0,0)=0$.

To retain the original zeta formula explicitly, the Mellin transform is
\begin{equation}\tag{DOR21}
\begin{gathered}
M_b(s)=2b^s\zeta(s)
\left\{\frac{\pi^2\Gamma((s+4)/2)}{2\pi^{(s+4)/2}}
-\frac{3\pi\Gamma((s+2)/2)}{4\pi^{(s+2)/2}}\right\},\\
M_b(0)=\frac14,\qquad M_b(1)=\frac b4,\qquad
M_{f_b^\sharp}(s)=\overline{M_b(1-\overline s)}.
\end{gathered}
\end{equation}
The first equality follows by absolute Mellin integration for
$\operatorname{Re}s>1$. The entire continuation and endpoint evaluations
are ETC10--19: both Gamma summands, the pole of the original zeta and
the endpoint zeros are retained there before evaluating their combined
limits. The sharp identity uses the original
$f^\sharp(u)=u^{-1}\overline{f(u^{-1})}$ and its measure $du/u$.

For prescribed complex-linear endpoint changes $(D,C)$, choose any
$b_1,b_2>0$ with $b_1\ne b_2$ and any fixed original finite vector
$k$. Set $a_i=N(k)b_i$ and retain the original ETC inverse:
\begin{equation}\tag{DOR22}
\begin{gathered}
c_1=\frac{4(D-b_2C)}{b_1-b_2},\qquad
c_2=\frac{4(b_1C-D)}{b_1-b_2},\qquad
t=c_1t_{k,a_1,y}+c_2t_{k,a_2,y}\in D\cap K_y^0,\\
(E\otimes E)t=0,\qquad (Q\otimes Q)t=0,\qquad
E\mu_y t=c_1f_{b_1}+c_2f_{b_2},\\
\bigl(M_{E\mu_y t}(1),M_{E\mu_y t}(0)\bigr)=(D,C).
\end{gathered}
\end{equation}
Substitution gives $b_1c_1+b_2c_2=4D$ and $c_1+c_2=4C$.
If codegree is defined instead as $M_{f^\sharp}(1)$, replace $C$
in the two coefficients by $\overline C$; DOR21 supplies exactly
the required conjugation. The original theta image $E\mu_y t$
belongs to $\mathcal V=E(T_0)$. Its action on the arithmetic
quotient is zero, and the original trace pairing with any admitted
second argument is zero by Connes--Consani--Marcolli,
Lemma \texttt{vanishV} and Proposition \texttt{pairTrvanishV}.
Those are cited arithmetic representation results, not new sign
theorems. DOR22 supplies their actual new tensor antecedent and every
one of its maps. Consequently both endpoint values can be varied
independently within the same arithmetic quotient class by this
explicit source family, while all original factors and their
trace contributions remain in DOR21 and the ETC comparison.

The lifted vanishing in DOR22 has the source's unchanged lattice label
as in DOR19. In particular these degree changes occur among
top-supported tensors and give supported zero under the individual
periodizations; neither the whole tensor nor its support is declared
absent. This computes the programme's exact relation with the CCM
trivial-correspondence operation without turning freedom to change
degree into an unproved positivity assertion.
